\section{Conclusions}
\label{sec:conc}

We have computed the quark--antiquark--gluon component of the outgoing state of a transversely
polarised virtual photon which scatters off a dense target at small $x$, for quarks of arbitrary mass,
and from it the leading-order trijet cross-section. The calculation uses light-cone perturbation theory
for the projectile and the eikonal approximation, in the CGC effective theory, for its interaction with
the target.

With respect to the massless case, the quark mass introduces helicity-flip terms in both vertices and
a second mass scale, $\of=\xi\mf/(1-\zeta)$, in the energy denominators. The amplitude is therefore a
sum over four helicity paths. For a gluon emitted before the collision each path is described by a
radial function with the one-parameter integral representation \eqref{eq:Gab}; for a gluon emitted
after the collision the radial functions factorise into a leading-order photon wave function and an
emission kernel off a massive fermion, eq.~\eqref{eq:Aexplicit}; and the instantaneous terms keep their
massless form with $Q\to\Ef$, eq.~\eqref{eq:GE}. The helicity and polarisation sums in the
cross-section, eqs.~\eqref{eq:H-direct}, \eqref{eq:Ht-direct} and \eqref{eq:Hx}, contain the massless
structures plus a small number of mass terms: the squares of the three paths with flips, the
interference of the no-flip path with the double-flip one and of the two single-flip paths with each
other, and the interference of the double-flip path with the instantaneous term. The colour structure
is the same as for massless quarks.

We have then studied three limits. For $\mf\to0$ we recover the results of ref.~\cite{Iancu:2022gpw}; the
limit is however not uniform in the collinear region and in the photoproduction region. For a soft
gluon the massive leading-order dijet cross-section factorises from the massless eikonal kernel, and we
recover the production JIMWLK Hamiltonian and the real part of the B--JIMWLK evolution of the
quadrupole; the same soft-gluon wave function gives the BK equation through the optical theorem. The
leading-logarithmic high-energy evolution is thus independent of the quark mass, which only enters the
impact factor. For a collinear gluon emitted after the collision the cross-section factorises into the
massive leading-order cross-section and the quasi-collinear splitting function of a massive quark,
eq.~\eqref{eq:Pmassive}, which displays the dead cone: the collinear singularity of the massless
calculation is replaced by a logarithm of the ratio of the hard scale to the quark mass.

Our results are exact in the kinematics of the projectile and resum the multiple scattering off the
target to all orders, but they are obtained at tree level and in the eikonal approximation, and for
the transverse polarisations of the photon only (section~\ref{sec:validity}). Several extensions
suggest themselves. The longitudinal photon should be added in order to describe electroproduction at
finite $Q^2$; the mass does not enter the longitudinal photon vertex, so that only two helicity paths
contribute. The virtual corrections with massive quarks are needed to complete the NLO calculation of
heavy-quark dijet production, and to check explicitly the cancellation of the soft divergences and the
emergence of the full B--JIMWLK equation; the one-loop photon wave functions of
refs.~\cite{Beuf:2021srj,Beuf:2022ndu} are the natural starting point. On the phenomenological side,
the coherent cross-section of section~\ref{sec:mom} can be evaluated numerically with a model for the
dipole $S$-matrix, and the complete cross-section can be used to study three-particle azimuthal
correlations of heavy-flavoured hadrons, in electroproduction at the Electron--Ion Collider and in
photoproduction in ultra-peripheral collisions. We leave these studies for future work.
